% Appendix B — compact half-page Computational Feasibility Pilot.
% LLM-ONLY pilot. Not a substitute for the IRB-gated human study.

\section*{Appendix B: Computational Feasibility Pilot}\label{sec:appendix-b}

Before the IRB-gated Phase~2 we ran an end-to-end pipeline pilot using
GPT-4o-mini both as generator and as LLM-as-judge rater. \PilotN{} problems
were sampled with a fixed seed from HumanEval~\cite{Chen2021HumanEval}
(164 Python tasks, MIT-licensed). Each problem produced nine conditions
(three CLEAR levels $\times$ three pedagogical framings: algorithm
explanation, complexity analysis, problem solving) for \PilotSessions{}
sessions. A Friedman test on the three prompt levels pooled across
problems yielded $\chi^2(\PilotFriedmanDF)=\PilotFriedmanChi$,
$p=\PilotFriedmanP$, Kendall's~$W=\PilotKendallW$; per-cell composite
means (basic / intermediate / advanced) were \PilotAlgoBasic{} /
\PilotAlgoInt{} / \PilotAlgoAdv{} for algorithm explanation,
\PilotCxBasic{} / \PilotCxInt{} / \PilotCxAdv{} for complexity analysis,
and \PilotPsBasic{} / \PilotPsInt{} / \PilotPsAdv{} for problem solving
on the 5--25 composite scale. Because LLM-as-judge is known to inflate
scores and compress between-level variance relative to calibrated human
raters, Kendall's~$W$ should be read as a \emph{floor} on the effect
size expected in the human study, not as a substitute for the
pre-registered partial $\eta^2\geq0.06$ target. Pilot data, prompts, the
rubric, and invocation logs accompany the supplementary material.
